\originalchapter{8.422公开作业一至四}
题源为Ketterle的Spring 2013公开作业; 第1、3份文件个别日期写2011, 与其课程页及到期年不一致, 原件版本信息在清单中保留. 本章均为编者解答, 对原题可检验的数学错误另作说明.
\originalsection{作业一: 光场状态}
\source{Assignment 1, 第1题.}
\begin{exercise}(a) 求任意两相干态重叠; (b) 证$\int\ket\alpha\bra\alpha\dd^2\alpha/\pi=I$; (c) 用中心对易子BCH公式证明$D(\alpha)=\exp(\alpha a^\dagger-\alpha^*a)$把真空移成相干态; (d) $E_x=E_0[a(t)+a^\dagger(t)]\sin kz$, 求自由相干态平均场与rms涨落, 解释为何涨落不依场强、时间, 真空也非零; (e) 给定相位分布$P(\phi)=|\sum_n\langle n|\alpha\rangle\e^{-\ii n\phi}|^2/(2\pi)$, 说明大$|\alpha|$的经典相位极限与数态无相位.\end{exercise}
\begin{solution}(a) 用数态级数求和得$\langle\alpha|\beta\rangle=\exp[-(|\alpha|^2+|\beta|^2)/2+\alpha^*\beta]$. (b) 在$m,n$矩阵元上, 角积分给$\delta_{mn}$, 径向积分$2\int_0^\infty r^{2n+1}\e^{-r^2}\dd r=n!$, 抵消归一化阶乘. (c) BCH得$D=\e^{-|\alpha|^2/2}\e^{\alpha a^\dagger}\e^{-\alpha^*a}$, 作用真空即相干态级数. (d) $\langle E_x\rangle=2E_0\Re(\alpha\e^{-\ii\omega t})\sin kz$, $\Delta E_x=|E_0\sin kz|$, 由$\langle aa^\dagger\rangle-\langle a^\dagger a\rangle=1$产生. 相干态只平移真空噪声; 场节点处此模式贡献为零. (e) 把$\alpha=|\alpha|\e^{\ii\theta}$代入, 所有级数相位在$\phi=\theta$同向, 大光子数时峰变窄, 相位标准差约$1/(2\sqrt{\bar n})$. 数态只有一个级数项, $P=1/(2\pi)$. 对任意态角积分的不同$n$交叉项为零, 也验证$P$归一化.
\end{solution}
\source{Assignment 1, 第2题.}
\begin{exercise}输入光态$\psi$和真空, 50:50分束器输出$a_1=(a+b)/\sqrt2,b_1=(a-b)/\sqrt2$. (a) 计数重合电压$V=V_0\langle n_{a_1}n_{b_1}\rangle$怎样测$g^{(2)}$; (b) 证经典$g^{(2)}(0)\ge1$; (c) 求相干态与$\ket2$的$g^{(2)}$; (d) 求$\ket\alpha\pm\ket{-\alpha}$两猫态的$g^{(2)}$, 是否非经典?\end{exercise}
\begin{solution}(a) 两输出模式对易, 乘积可正规序重排; 真空项消失, 得$V/V_0=\langle a^{\dagger2}a^2\rangle/4$. 单路均值为$\bar n/2$, 故除两单路平均的乘积即$g^{(2)}$, 理想模型无加性偏移, 实际电子背景可另校正. 原题单模定常式不含非平凡$\tau$依赖. (b) 非负强度方差给结论. (c) 分别为1和$1/2$. (d) 原题$1/\sqrt2$未归一化有限重叠的猫态, 应用$N_\pm=\sqrt{2(1\pm\e^{-2x})}$,$x=|\alpha|^2$. 平均数$x\tanh x$与$x\coth x$, 二阶阶乘矩均为$x^2$, 故$g_+^{(2)}=\coth^2x>1$,$g_-^{(2)}=\tanh^2x<1$. 奇猫表现亚Poisson非经典相关; 偶猫虽无此证据仍可为非经典态, 不能从$g^{(2)}>1$推断经典.
\end{solution}
\source{Assignment 1, 第3题.}
\begin{exercise}题定义$Q_\rho(z)=\langle z|\rho|z\rangle$, 不含$1/\pi$. 求并绘: (i) $\cos(\theta/2)\ket0+\e^{\ii\phi}\sin(\theta/2)\ket{12}$, 多组$\theta,\phi$, 是否最小不确定或压缩; (ii) 偶猫$\alpha=3$, 看对数$Q$干涉; (iii) $N^{-1/2}\sum_{k=1}^N\e^{\ii k\phi}\ket k$, $N=10,\phi=\pi/4$, 解释$N\to\infty$及压缩; (iv) 压缩真空$r=0.2,1.2,4$; (v) Kerr演化$\e^{\ii\xi n(n-1)t}\ket\beta$, $\beta=4,\xi=\pi/128$, 多时刻图, 求双猫时间与复现时间.\end{exercise}
\begin{solution}(i) $Q=\e^{-|z|^2}|\cos(\theta/2)+\e^{\ii\phi}\sin(\theta/2)(z^*)^{12}/\sqrt{12!}|^2$. $\langle a\rangle=\langle a^2\rangle=0$, 两正交分量方差均$[1+24\sin^2(\theta/2)]/4$, 除真空端点外非最小且不压缩. (ii) $Q=\e^{-|z|^2-9}|\e^{3z^*}+\e^{-3z^*}|^2/[2(1+\e^{-18})]$, 对数图显出相干和的干涉, 与两相干态混合的$Q$不同. (iii) $Q=\e^{-|z|^2}|\sum_{k=1}^N\e^{\ii k\phi}(z^*)^k/\sqrt{k!}|^2/N$. 相位级数相干集中到$\phi$方向, 光子数均匀分布在1至$N$. $N\to\infty$不收敛成归一化有限能量态, 平均数发散. 它是相位局域化的非Gaussian态, 不可一概称正交分量压缩; 实际二阶矩检查下$N=10$最小正交分量方差约0.610, 大于真空$1/4$. (iv) 按题面正系数的压缩约定,$Q=\e^{-|z|^2}\exp[\tanh r\Re(z^{*2})]/\cosh r$. $Q$包含测量真空噪声, 椭圆窄轴不能无限缩至零, 不能用$Q$轮廓宽度直接当$X$方差. (v) $Q=\e^{-|z|^2-|\beta|^2}|\sum_n(\beta z^*)^n\e^{\ii\xi n(n-1)t}/n!|^2$. $t=64$时相位序列可分解成$\ket{\ii\beta}$与$\ket{-\ii\beta}$的等权相干和; $t=128$时所有$n(n-1)$为偶数, 完全复现. 绘图与截断误差核验见随卷的自绘图与数值脚本.
\end{solution}

\originalsection{作业二: 分束器与纠缠}
\source{Assignment 2, 第1题.}
\begin{exercise}$B=\exp[\theta(a^\dagger b-ab^\dagger)]$. (a) 证光子总数守恒与幺正; (b) 求$B\ket\alpha_a\ket0_b$, 给透射反射率; (c) 定义$s_z=n_a-n_b,s_+=a^\dagger b,s_-=ab^\dagger,s_\pm=(s_x\pm\ii s_y)/2$, 写$B$的自旋形式, 证对易关系, 说明总数与Casimir的关系; (d) 求50:50下单端数态$\ket n$的输出和$n=10,100$分布宽度.\end{exercise}
\begin{solution}(a) 生成元反Hermitian且与$n_a+n_b$对易. (b) 以题面指数$G=a^\dagger b-ab^\dagger$为定义, $[G,a]=-b,[G,b]=a$, 因而
\[BaB^\dagger=a\cos\theta-b\sin\theta,\quad
 BbB^\dagger=a\sin\theta+b\cos\theta.\]
原件后给的式(4)把这一对变换的符号写成反向, 与所列指数不一致; 本解和第3题(d)统一按指数. 对产生算符作同样变换, 得$B\ket\alpha\ket0=\ket{\alpha\cos\theta}\ket{-\alpha\sin\theta}$, 强度率仍$\cos^2\theta,\sin^2\theta$. (c) $B=\exp(\ii\theta s_y)$,$[s_i,s_j]=2\ii\epsilon_{ijk}s_k$. 固定总数$N$的表示自旋$j=N/2$, $s_x^2+s_y^2+s_z^2=N(N+2)I$. 总数与Casimir相关, 但总数本身不是字面上的二次Casimir, 原题表述需此修正. (d) 按同一指数约定, 输出$\sum_{k=0}^n(-1)^{n-k}\sqrt{\binom nk}2^{-n/2}\ket{k,n-k}$. 单路平均$n/2$, 标准差$\sqrt n/2$, $n=10,100$分别约1.58,5, 相对标准差$1/\sqrt n$. 最大概率在最接近$n/2$处, 不形成精确的一半一半数态.
\end{solution}
\source{Assignment 2, 第2题.}
\begin{exercise}(a) 错误衰减模型$a\mapsto ta$, $|t|<1$, 为何使相干态光子数均值$|t|^2|\alpha|^2$而方差$|t|^4|\alpha|^2$, 产生不物理结果? (b) 改用$a_{\rm out}=ta+rb$,$b$真空,$|t|^2+|r|^2=1$, 证保持散粒噪声及相干态输出.\end{exercise}
\begin{solution}(a) 它把$n$错误地当作连续缩放的随机量, 又使$[a_{\rm out},a_{\rm out}^\dagger]=|t|^2\ne1$, 不能是单模式幺正变换. 二正交分量的不确定积也被缩小. 亚Poisson本身并非违背不确定关系, 错在这一映射不保持量子代数. (b) 加真空端口后对易子恢复1, 输出相干态$\ket{t\alpha}$, 均值与方差同为$|t\alpha|^2$. 一般输入方差变为$\eta^2\operatorname{Var}n+\eta(1-\eta)\bar n$, 后项为损失随机性的附加噪声.
\end{solution}
\source{Assignment 2, 第3题.}
\begin{exercise}纯态Schmidt数为非零系数个数. (a) 证为1当且仅当直积; (b) 证局域幺正不改变它, 检查原题随后声称任意LOCC不改变它是否成立; (c) 求$(00+11+22)/\sqrt3$, $(00+01+10+11)/2$, $(00+01+10-11)/2$, $(00+01+11)/\sqrt3$的数; (d) Yurke态$(\ket{n,n-1}+\ket{n-1,n})/\sqrt2$输入及经本题实50:50分束器后的Schmidt数, 相干加真空的数; (e) 原件同时写$S_{\rm op}(r)=\exp[-(r/2)(ab-a^\dagger b^\dagger)]$与系数$\lambda_k=\operatorname{sech}r\tanh^kr$. 检查这两种$r$约定, 再按给定系数绘大于0.01的个数, 检查其随$r$必单调增长的判断.\end{exercise}
\begin{solution}(a) 系数矩阵的SVD秩为1当且仅当为外积. (b) 局域幺正左右乘不改奇异值, 但局域测量可把纠缠态投影成直积, 所以任意LOCC“不改变”是错的. 对单个成功分支, 局域线性算符不能增加矩阵秩; 不增加与不改变不同. (c) 依次3,1,2,2, 后两者二乘二系数矩阵行列式非零. (d) 输入为2. 用生成多项式, 按第1题指数, 输出正比$(a^{\dagger2}-b^{\dagger2})^{n-1}a^\dagger\ket0$, 有$n$个非零固定总数项, 所以为$n$, $n=1$时反而为1. 相干输入输出直积, 为1. 分束器对两端分割是非局域操作, 不与(b)矛盾; Schmidt秩并非相位灵敏度的充分定量指标. (e) 生成元的系数决定物理压缩量. 若采用原件指数, 真正归一化系数为$\operatorname{sech}(r/2)\tanh^k(r/2)$; 原件列的$\operatorname{sech}r\tanh^kr$对应$\exp[r(a^\dagger b^\dagger-ab)]$. 以下及阈值图按题目所列系数定义$r$, 若按指数横轴则将$r$加倍. 对$0<\tanh r<1$, 满足$\lambda_k>\epsilon$的整数$k\ge0$为$k<\ln(\epsilon\cosh r)/\ln(\tanh r)$. 计数由这个严格不等式取整, 若$\operatorname{sech}r\le\epsilon$则为0. $r\ge\operatorname{arcosh}100$时全部系数都不大于0.01（边界$\lambda_0=0.01$）, 所以严格大于的计数为0; 数值$\operatorname{arcosh}100\simeq5.298292$仅作近似显示, 故原题“单调增长”不成立. 真实纠缠熵随$r$增长, 固定幅度阈值的计数却在大$r$失去意义.
\end{solution}

\originalsection{作业三: 损耗与压缩}
\source{Assignment 3, 第1题.}
\begin{exercise}Mach--Zehnder单光子输入, 理想$P_b=\cos^2(\phi/2)$, 可见度$V=(P_{\max}-P_{\min})/(P_{\max}+P_{\min})$. (a) 无损$V$; (b) 一臂与真空损失口混合角$\theta$, 写各步状态、概率、可见度; (c) 两臂角$\theta,\theta'$, 同样计算, 只保留存活光子时可见度怎样变, 是否只取决于$\Delta\theta$? (d) 相干态输入的可见度是否变化, 哪种统计能区分单光子和相干态?\end{exercise}
\begin{solution}(a) $V=1$. 选相位约定使第一分束器后态$(\ket a+\ket b)/\sqrt2$. (b) 经相移与损失后态$[\ket a+t\e^{\ii\phi}\ket b+r\e^{\ii\phi}\ket c]/\sqrt2$, $t=\cos\theta,r=\sin\theta$. 最后分束器后存活两端振幅$(1\pm t\e^{\ii\phi})/2$, 损失振幅$r\e^{\ii\phi}/\sqrt2$. $P_b=(1+t^2+2t\cos\phi)/4$, $V=2t/(1+t^2)$. (c) 两损失口$c,d$必须分别保留, 中间态$[t'\ket a+r'\ket d+\e^{\ii\phi}(t\ket b+r\ket c)]/\sqrt2$, 末态存活振幅$(t'\pm t\e^{\ii\phi})/2$, $P_b=(t'^2+t^2+2tt'\cos\phi)/4$, $V=2tt'/(t'^2+t^2)$. 一般不只由$\Delta\theta$决定; 平衡$t=t'$时为1. 条件存活概率$(t'^2+t^2)/2$与$\phi$无关, 归一化不改可见度. (d) 线性装置对相干幅度给相同平均强度可见度, 但单光子有总计数至多1、两端反相关; 相干态有Poisson总数和非零双计数概率. 平均条纹不能区分这些统计.
\end{solution}
\source{Assignment 3, 第2题.}
\begin{exercise}$H/\hbar=\omega a^\dagger a+\ii\Lambda[(a^\dagger)^2\e^{-2\ii\omega t}-a^2\e^{2\ii\omega t}]$, $\Lambda$实. (a) 取$a=b\e^{-\ii\omega t}$写$b,b^\dagger$方程; (b) $P=(b+b^\dagger)/2,Q=(b-b^\dagger)/(2\ii)$怎样对应场分量, 求演化; (c) 真空初态的光子数与两标准差.\end{exercise}
\begin{solution}$\dot a=-\ii\omega a+2\Lambda a^\dagger\e^{-2\ii\omega t}$, 故$\dot b=2\Lambda b^\dagger,\dot b^\dagger=2\Lambda b$. 所以$\dot P=2\Lambda P,\dot Q=-2\Lambda Q$, $P(t)=\e^{2\Lambda t}P(0),Q(t)=\e^{-2\Lambda t}Q(0)$, $b(t)=b(0)\cosh r+b^\dagger(0)\sinh r$, $r=2\Lambda t$. 场为这两算符乘空间时间正交的正弦余弦项. 真空$\bar n=\sinh^2r$, $\Delta P=\e^r/2,\Delta Q=\e^{-r}/2$, 不确定积保持$1/4$, 泵浦向场输入光子对能量.
\end{solution}
\source{Assignment 3, 第3题.}
\begin{exercise}压缩真空$r$, 再经透过率$T$的损失口. (a) 光子数方差、与分量方差关系、大$r$标度及相干态比较; (b) 定义压缩质量为大/小分量方差比, 求它降为原来一半的$T$; (c) 损失口改用沿窄轴的相干态LO,$\alpha$, 求输出相空间、平均光子数、光子数方差各贡献及大$r$小损耗极限, 给亚Poisson条件并解释.\end{exercise}
\begin{solution}(a) Gaussian Wick收缩给$\operatorname{Var}n=2\bar n(\bar n+1)=\sinh^2(2r)/2$, 其中$\bar n=V_++V_--1/2$, $V_\pm=\e^{\pm2r}/4$. 大$r$方差约$2\bar n^2$, 大于相干态的$\bar n$. (b) 新方差比$[1+T(\e^{2r}-1)]/[1+T(\e^{-2r}-1)]$. 令它等于$\e^{4r}/2$得$T=(\e^{4r}-2)/[\e^{4r}+\e^{2r}-2]$. 若$\e^{4r}<2$, 目标低于无压缩的1, 物理$T$区间无解. 大$r$允许损失约$\e^{-2r}$.

(c) 输出为衰减椭圆加位移$\beta=\sqrt{1-T}\alpha$, 宽窄轴方差如(b). $N=T\sinh^2r+|\beta|^2$. 令$n_s=\sinh^2r$, 无位移部分方差$2T^2n_s(n_s+1)+T(1-T)n_s$. 加位移后多$|\beta|^2[1-T+T\e^{-2r}]$, 奇数阶混合矩因压缩真空偶宇称为零. 亚Poisson要求
\[|\beta|^2>\frac{T(2n_s^2+n_s)}{1-\e^{-2r}},\qquad |\beta|^2=(1-T)|\alpha|^2.\]
小损耗且大$r$时阈值约$2n_s^2$. 强LO把窄分量涨落映成主要光子数涨落, 但还须压过原压缩真空大的二阶光子对噪声. 不能只说“有压缩就一定亚Poisson”.
\end{solution}
\source{Assignment 3, 第4题.}
\begin{exercise}Alice有未知$a\ket0+b\ket1$及EPR对$(00+11)/\sqrt2$的第一端, Bob有第二端. Alice测$\Phi_\pm=(00\pm11)/\sqrt2,\Psi_\pm=(01\pm10)/\sqrt2$, 各结果Bob态是什么, 怎样用两经典比特恢复? 原件随后连续变量部分为阅读背景, 不设额外必答小问.\end{exercise}
\begin{solution}把三比特重写成Bell基, 得$\frac12[\ket{\Phi_+}\ket\psi+\ket{\Phi_-}Z\ket\psi+\ket{\Psi_+}X\ket\psi+\ket{\Psi_-}XZ\ket\psi]$, 最后一项整体相位无关. 每结果概率$1/4$, Bob据两个比特做逆Pauli. 未得消息前Bob态为$I/2$, 所以没有超光速信息. 连续变量的分束器与同相测量可实现联合正交分量测量, 有限压缩带来噪声; 原件无限压缩态非正规化, 不能当作有限能量实物.
\end{solution}

\originalsection{作业四: Yurke与原子关联}
\source{Assignment 4, 第1题.}
\begin{exercise}Yurke态$(\ket{n,n-1}+\ket{n-1,n})/\sqrt2$, 分束器$BaB^\dagger=(a+\ii b)/\sqrt2$, $BbB^\dagger=(b+\ii a)/\sqrt2$, 相移$PaP^\dagger=a\e^{\ii\phi}$, 装置$B^\dagger PB$. (a) 给三个中间态, 检查零相位复现; (b) 测$M=n_a-n_b$, 在$\phi=0$算误差传播, 原题期望$\sqrt2/n$; (c,d) 相干态$\ket\alpha\ket0$的各步状态与相位误差.\end{exercise}
\begin{solution}(a) 令$F(x,y)=[x^ny^{n-1}+x^{n-1}y^n]/\sqrt{2n!(n-1)!}$. 第一段后为$F[(a^\dagger-\ii b^\dagger)/\sqrt2,(b^\dagger-\ii a^\dagger)/\sqrt2]\ket0$. 相移后明确为
\[F[(\e^{-\ii\phi}a^\dagger-\ii b^\dagger)/\sqrt2,\,
 (b^\dagger-\ii\e^{-\ii\phi}a^\dagger)/\sqrt2]\ket0.\]
第三段后为下列输出. 等价输出为$F[a^\dagger\e^{-\ii\phi/2}\cos(\phi/2)+b^\dagger\e^{-\ii\phi/2}\sin(\phi/2),\ b^\dagger\e^{-\ii\phi/2}\cos(\phi/2)-a^\dagger\e^{-\ii\phi/2}\sin(\phi/2)]\ket0$, 上述符号依本题固定B/P定义, $\phi=0$即原态.

(b) 固定总数$2n-1$, $J_z=M/2$初态仅含$m=\pm1/2$. $\langle M\rangle=0,\operatorname{Var}M=1$, $\langle J_x\rangle=n/2$. 装置使$M(\phi)=2(J_z\cos\phi\pm J_x\sin\phi)$, 零点斜率绝对值$n$, 故$\delta\phi=1/n$. 按题给态、装置和误差公式不能推出$\sqrt2/n$; 原件目标系数有误, 本卷保留可复算结果. (c) 相干幅度依次$(\alpha/\sqrt2,-\ii\alpha/\sqrt2)$、$(\alpha\e^{-\ii\phi}/\sqrt2,-\ii\alpha/\sqrt2)$、$(\alpha\e^{-\ii\phi/2}\cos(\phi/2),+\alpha\e^{-\ii\phi/2}\sin(\phi/2))$. 第二分量正号由本题固定B/P直接相乘确定. (d) 输出数差平均$\bar n\cos\phi$, 方差$\bar n$, 所以$\delta\phi=1/(\sqrt{\bar n}|\sin\phi|)$. 最佳在$\pi/2$, 零点斜率为零不能用同一测量局域估计.
\end{solution}
\source{Assignment 4, 第2题.}
\begin{exercise}忽略重力, 两粒子自源点$A,B$到探测点1,2的振幅按$\e^{\ii\mathbf k\cdot\mathbf r}$传播. (a) 求交换对称/反对称的联合概率对$\mathbf r_{21}$的依赖; (b) 源横宽$W$, 探测宽$w$, 距离$d\gg W,w$, 证需$Ww\ll d\lambda_{\rm dB}$; $^6$Li500$\mu$K,$d=10$ cm, $W\simeq w$, 估波长和尺寸上限; (c) 波矢差一维Gaussian$\propto\e^{-q^2/\gamma^2}$, 求平均关联和宽度; 脉冲源纵宽$L$, 飞行时间$\tau$证明$\Delta k_l\lesssim mvL/(\hbar d)$, 用$\tau=0.1$s,$L\simeq W$估定时分辨率; (d) 证相干单元空间体积增加$d^3/(W^2L)$, 估数量级与平均相空间占据, 比较BEC/Fermi退化气.\end{exercise}
\begin{solution}(a) 两路径相位差为$(\mathbf k_A-\mathbf k_B)\cdot\mathbf r_{21}$, 概率$2|\psi_A\psi_B|^2[1\pm\cos(\Delta\mathbf k\cdot\mathbf r_{21})]$. (b) 横向角差$W/d$, 相位量级$2\pi Ww/(d\lambda)$, 要远小于$2\pi$. 采用一维热动量尺度$\sqrt{mk_BT}$,$\lambda=h/\sqrt{mk_BT}\simeq80$ nm; 用三维rms速度则为46 nm, 必须说明约定. $W\ll\sqrt{d\lambda}\simeq89\,\mu$m. (c) 归一化Gaussian特征函数给$\overline P=2|\psi_A\psi_B|^2[1\pm\e^{-\gamma^2r_l^2/4}]$, 宽度尺度$2/\gamma$. 同一探测时刻初始纵向差$L$只允许速度差$L/\tau=vL/d$, 即题给界. 相干纵长约$h/\Delta p\gtrsim\lambda d/L$, 时间分辨率需小于$(\lambda d/L)/v=\lambda\tau/L$. 以尺寸边界89$\mu$m估约90$\mu$s, 取更小源时容许更长相干时间.

(d) 自由线性膨胀使三方向每个相干长度分别增加$d/W,d/W,d/L$, 相乘即题给因子. 边界尺寸下约$1.4\times10^9$, 实际严格准直需$W$更小. Liouville定理使相空间密度保持, 平均占据为$n_{\rm atoms}\lambda_T^3$, 题目未给原子数或密度, 所以不能数值确定; 对常见MOT通常远小于1是额外实验经验, 非题给条件的推论. BEC有宏观最低模式占据, 退化Fermi每内部态每单粒子模式占据不超过1.
\end{solution}
